\section{Implementation}\label{implementation}

This chapter reports the technology-dependent choices made while turning the
design of \cref{design} into a running system, following the same structure
(network protocols, data representation, database access, authentication and
authorisation) before detailing the concrete frameworks used.

\paragraph{Network protocols.} Three protocols are used, matching the three
communication styles of \cref{architecture}. \textbf{HTTP/1.1} carries every
synchronous REST call, both client-to-gateway and gateway-to-service, and is
also the transport HTTP/1.1 upgrades from for \textbf{WebSocket}: every
Socket.IO server (CAD collaboration, notifications, chatbot,
\cref{modelling}) is configured with \texttt{transports:
['polling','websocket']}, so a client starts with HTTP long-polling and
transparently upgrades to a persistent WebSocket connection where the network
allows it, falling back gracefully otherwise (e.g.\ behind a restrictive
proxy). Internally, the \textbf{Redis protocol (RESP)} carries pub/sub domain
events between the catalog, cart and notification services via
\texttt{ioredis} clients (\cref{interaction}). Finally, the chatbot service is
the only component that talks to the outside world beyond the browser: it
issues outbound \textbf{HTTPS} requests to the Mistral AI chat-completions
API (\texttt{https://api.mistral.ai/v1/chat/completions}) to generate
answers, with a template-based fallback if no API key is configured.

\paragraph{In-transit data representation.} \textbf{JSON} is the only
wire format used anywhere --- REST request/response bodies, Redis event
payloads, and every Socket.IO message --- keeping the whole system
technologically homogeneous and human-readable for debugging. A few
conventions are applied consistently across services to avoid ambiguity:
identifiers are \textbf{string UUIDs} used directly as MongoDB \texttt{\_id}
(not the default \texttt{ObjectId}), so an id looks and behaves the same
whether logged, sent over REST, or embedded in a WebSocket payload; monetary
amounts are represented as \textbf{integer euro cents} rather than
floating-point decimals, both in the catalog's source price and everywhere it
is copied (cart line items, order snapshots, BOM items), specifically to
avoid floating-point rounding errors in a domain full of price sums; and
timestamps are serialised as \textbf{ISO~8601} strings at the HTTP boundary.

\paragraph{Database access.} Every service talks to its own MongoDB
exclusively through \textbf{Mongoose} (an ODM), never the raw driver: domain
entities are mapped to schemas in each service's
\texttt{adapters/persistence} layer, and repositories expose only
domain-shaped methods to the rest of the service (a port/adapter boundary,
\cref{technological-details}). Queries are plain Mongoose finds/filters with
indexes (\cref{data-and-consistency-issues}) for the transactional read/write
paths; the one place using MongoDB's \textbf{aggregation pipeline} is the
reporting service (\texttt{MongoOrderTrendReader}), which groups orders by
day and status (\texttt{\$match}/\texttt{\$group}/\texttt{\$sum}) directly in
the database rather than pulling raw documents and summing them in
application code. Optimistic concurrency (\cref{data-and-consistency-issues})
is implemented with a conditional \texttt{findOneAndUpdate} rather than
MongoDB transactions, consistent with the design decision to avoid
multi-document transactions altogether.

\paragraph{Authentication.} Implemented with \textbf{JWT} (the
\texttt{jsonwebtoken} library, HS256), issued by the authentication service on
login/registration/guest-session creation and verified exclusively at the API
gateway (\cref{security}); passwords are hashed with \textbf{bcrypt}
(\texttt{bcryptjs}, 12 salt rounds) and never stored or logged in clear text.

\paragraph{Authorisation.} Implemented as plain \textbf{RBAC} via a small
hand-written Express middleware duplicated in each service that needs it (no
shared authorisation package), reading the gateway-issued
\texttt{X-User-Role} header, complemented by ownership checks written
directly in use-cases (comparing a resource's \texttt{userId}/\texttt{ownerId}
against the caller's identity) rather than a generic ABAC/policy engine ---
adequate given the small, fixed role set (\texttt{GUEST}/\texttt{BASE}/
\texttt{ADMIN}) and the fact that most access rules are simple ownership
checks rather than complex attribute combinations.

\subsection{Technological details}\label{technological-details}

\paragraph{Backend.} All nine backend components (the gateway and the eight
domain services) are written in \textbf{TypeScript~5.4} on \textbf{Node.js~24}
and follow the same internal layering (a light hexagonal/clean-architecture
style): a \texttt{domain} layer (entities, value objects, port interfaces)
with no framework dependency, a \texttt{useCases} layer implementing
application logic against those ports, and an \texttt{adapters} layer
(\texttt{http}, \texttt{persistence}, \texttt{httpClient}, \texttt{messaging},
\texttt{websocket} as needed) wiring the use-cases to \textbf{Express~4.18},
\textbf{Mongoose~8.4}, \textbf{ioredis} and \textbf{socket.io}. This uniform
structure is what let unit tests fake every port and never depend on a real
MongoDB (\cref{validation}). Cross-cutting libraries: \texttt{jsonwebtoken}
and \texttt{bcryptjs} (authentication service), \texttt{helmet} and
\texttt{cors} (gateway), \texttt{http-proxy-middleware} (gateway's REST/WS
proxying).

\paragraph{Frontend.} A \textbf{Vue~3.4} Single Page Application, written
entirely with the \textbf{Composition API} (\texttt{<script setup lang="ts">}),
bundled with \textbf{Vite~5}, and type-checked in strict mode with
\texttt{vue-tsc}. State is managed with \textbf{Pinia} (a store per
cross-cutting concern: \texttt{authStore}, \texttt{cartStore},
\texttt{notificationStore}), routing with \textbf{vue-router} (a single
route table with per-route \texttt{meta} flags for auth/role guards,
\cref{user-guide}). HTTP calls go through a small hand-written
\texttt{fetch}-based client rather than a library such as axios, and
real-time channels use \textbf{socket.io-client} through three dedicated
wrappers, one per Socket.IO server (CAD collaboration, notifications,
chatbot), mirroring the backend's separation. No CSS framework is used ---
components are styled with scoped \texttt{<style>} blocks over a small shared
design-token stylesheet --- and, unlike the backend, the frontend currently
has no automated test suite (\cref{validation}).

\paragraph{External integration.} The chatbot is the only feature depending
on a third-party service: \textbf{Mistral AI}'s chat-completions endpoint
(default model \texttt{mistral-small-latest}), called with a bounded retry
and timeout policy (\cref{fault-tolerance}) and gated by an API key read from
an environment variable / Kubernetes secret --- never hard-coded in source.

\paragraph{Containerisation.} Every backend service and the frontend are
packaged as Docker images built \texttt{FROM node:24-alpine}; the frontend's
production image is a two-stage build that compiles the Vite bundle and then
serves it through a lightweight \textbf{nginx:alpine} image, which also acts
as the reverse proxy for \texttt{/api} traffic (REST and WebSocket upgrades
alike) as described in \cref{infrastructure} and detailed in
\cref{deployment}.
